% 20260916 Science Quiz mathematics Problem set.

\begin{problem}{1}
$U,V$가 $\R^{10}$의 subspaces이고
\[
    \dim U=6,
    \qquad
    \dim V=7
\]
이라 하자. $\dim(U\cap V)$가 가질 수 있는 최솟값을 구하시오.
\end{problem}

\begin{solution}
subspace dimension formula에 의해
\[
    \dim(U+V)
    =\dim U+\dim V-\dim(U\cap V).
\]
한편 $U+V\subseteq\R^{10}$이므로
\[
    \dim(U+V)\leq10.
\]
따라서
\[
    10
    \geq6+7-\dim(U\cap V),
\]
즉
\[
    \dim(U\cap V)\geq3.
\]
이 bound는 실제로 attain된다. 예를 들어
\[
    U=\Span(\mathbf e_1,\ldots,\mathbf e_6),
    \qquad
    V=\Span(\mathbf e_4,\ldots,\mathbf e_{10})
\]
로 두면
\[
    U\cap V=\Span(\mathbf e_4,\mathbf e_5,\mathbf e_6)
\]
이므로 dimension이 $3$이다. 따라서 최솟값은
\[
    \boxed{3}.
\]
핵심은 두 subspace의 dimension 합이 ambient dimension을 넘으면 그 초과분만큼 intersection이 강제된다는 것이다.
\end{solution}

% Current handbook anchor: 2026 Fields Medal, John Pardon.
\begin{problem}{2}
2026 Fields Medalist 중 symplectic geometry와 topology, holomorphic curves, virtual fundamental cycles, Fukaya categories에 대한 업적으로 알려진 수학자의 이름을 쓰시오.
\end{problem}

\begin{solution}
2026 Fields Medalist 네 명 가운데 symplectic geometry와 topology를 주된 분야로 하고, holomorphic curves, virtual fundamental cycles, Fukaya categories 등에 대한 업적으로 알려진 수학자는 John Pardon이다. 따라서 정답은
\[
    \boxed{\text{John Pardon}}.
\]
2026 Fields Medal의 짧은 association으로는
\[
    \text{John Pardon}
    \quad\longleftrightarrow\quad
    \text{symplectic geometry and topology}
\]
를 기억하면 좋다.
\end{solution}

\begin{problem}{3}
함수
\[
    f(x)=x^3+x+1
\]
의 inverse function을 $g=f^{-1}$라 하자. $g'(3)$을 구하시오.
\end{problem}

\begin{solution}
먼저
\[
    f(1)=1+1+1=3
\]
이므로
\[
    g(3)=1.
\]
Derivative of an inverse function에 의해
\[
    g'(3)
    =\frac{1}{f'(g(3))}
    =\frac{1}{f'(1)}.
\]
그런데
\[
    f'(x)=3x^2+1
\]
이므로
\[
    f'(1)=4.
\]
따라서
\[
    \boxed{g'(3)=\frac14}.
\]
Inverse function의 derivative를 물으면 inverse를 직접 구하려 하지 말고 원래 함수에서 대응하는 point를 먼저 찾는 것이 핵심이다.
\end{solution}

\begin{problem}{4}
집합 $\{1,2,\ldots,6\}$의 permutation을 uniformly random하게 하나 고른다. 숫자 $1,2,3$이 순서에 관계없이 서로 이웃한 세 자리 전체를 차지할 확률을 구하시오.
\end{problem}

\begin{solution}
$1,2,3$이 서로 이웃한 세 자리를 차지한다면 이 세 수를 하나의 block으로 볼 수 있다. Block 내부에서는
\[
    3!
\]
개의 순서가 가능하고, 이 block과 $4,5,6$을 합친 네 objects의 배열은
\[
    4!
\]
개이다. 따라서 favorable permutations의 수는
\[
    4!\,3!
\]
이고, 전체 permutation은 $6!$개이므로 확률은
\[
    \frac{4!\,3!}{6!}
    =\boxed{\frac15}.
\]
여러 지정된 objects가 consecutive해야 하는 counting에서는 그 objects를 하나의 block으로 묶는 것이 표준적인 첫 단계이다.
\end{solution}

% Current handbook anchor: POSTECH professor, KAIST Ph.D., 2025 Korea Science Award.
\begin{problem}{5}
2025 Korea Science Award를 수상한 \POSTECHKorean{} 교수로, \textup{KAIST}에서 Ph.D.를 받았으며 geometric topology, 특히 four-dimensional topology에 관한 연구로 알려진 수학자의 이름을 쓰시오.
\end{problem}

\begin{solution}
최종 본교재의 2026 Republic of Korea section에서 이 association에 해당하는 수학자는 Jae Choon Cha이다. 그는 \POSTECHKorean{} 교수이며 KAIST에서 Ph.D.를 받았고, geometric topology에서의 연구로 2025 Korea Science Award를 수상했다. 따라서 정답은
\[
    \boxed{\text{Jae Choon Cha}}.
\]
Science Quiz에서는
\[
    \text{POSTECH professor}
    \quad+\quad
    \text{KAIST Ph.D.}
    \quad+\quad
    \text{geometric topology}
\]
라는 양교 연결을 함께 기억해 두면 유용하다.
\end{solution}

\begin{problem}{6}
다음 정적분의 값을 구하시오.
\[
    \int_1^3
    \frac{1}{\sqrt{x+1}-\sqrt{x-1}}\,\dd x.
\]
\end{problem}

\begin{solution}
분모의 conjugate를 이용하면
\[
    \frac{1}{\sqrt{x+1}-\sqrt{x-1}}
    =\frac{\sqrt{x+1}+\sqrt{x-1}}{(x+1)-(x-1)}
    =\frac{\sqrt{x+1}+\sqrt{x-1}}{2}.
\]
따라서
\[
    \begin{aligned}
    \int_1^3
    \frac{1}{\sqrt{x+1}-\sqrt{x-1}}\,\dd x
    &=%
    \frac12\int_1^3
    \left(\sqrt{x+1}+\sqrt{x-1}\right)\,\dd x\\
    &=\frac13
    \left[(x+1)^{3/2}+(x-1)^{3/2}\right]_1^3\\
    &=\frac13\left(8+2\sqrt2-2\sqrt2\right)\\
    &=\boxed{\frac83}.
    \end{aligned}
\]
Radical이 분모에서 차 형태로 나타나면 substitution보다 먼저 conjugate rationalization을 확인하는 것이 좋다.
\end{solution}

\begin{problem}{7}
positive divisor가 정확히 $12$개인 가장 작은 positive integer를 구하시오.
\end{problem}

\begin{solution}
positive integer의 prime factorization을
\[
    n=p_1^{\alpha_1}\cdots p_k^{\alpha_k}
\]
라 하면 positive divisor의 개수는
\[
    \tau(n)=(\alpha_1+1)\cdots(\alpha_k+1)
\]
이다. 따라서 $\tau(n)=12$가 되려면 exponent tuple은
\[
    (11),
    \qquad
    (5,1),
    \qquad
    (3,2),
    \qquad
    (2,1,1)
\]
중 하나이다. 각 경우에서 가장 작은 integer는 각각
\[
    2^{11},
    \qquad
    2^5\cdot3,
    \qquad
    2^3\cdot3^2,
    \qquad
    2^2\cdot3\cdot5.
\]
이 가운데 가장 작은 것은
\[
    2^2\cdot3\cdot5=60.
\]
실제로
\[
    \tau(60)=(2+1)(1+1)(1+1)=12.
\]
따라서 정답은
\[
    \boxed{60}.
\]
주어진 divisor count를 factorize하여 exponent pattern으로 바꾸는 것이 핵심이다.
\end{solution}

% Current handbook anchor: Cleve Moler, died May 20, 2026.
\begin{problem}{8}
2026년에 별세한 numerical computing의 pioneer로, MATLAB을 만들고 MathWorks를 공동 설립한 인물의 이름을 쓰시오.
\end{problem}

\begin{solution}
최종 본교재의 2026 mathematicians and institutions section에서 2026년 5월 20일 별세한 Cleve Moler는 numerical computing의 pioneer이자 MATLAB의 creator, MathWorks의 cofounder로 정리되어 있다. 따라서 정답은
\[
    \boxed{\text{Cleve Moler}}.
\]
가장 짧은 quiz association은
\[
    \text{Cleve Moler}
    \quad\longleftrightarrow\quad
    \text{MATLAB and numerical computing}
\]
이다.
\end{solution}

\begin{problem}{9}
함수 $y$가 initial value problem
\[
    y''-3y'+2y=0,
    \qquad
    y(0)=0,
    \qquad
    y'(0)=1
\]
을 만족한다. $y(\ln 2)$를 구하시오.
\end{problem}

\begin{solution}
characteristic equation은
\[
    r^2-3r+2=0,
\]
즉
\[
    (r-1)(r-2)=0
\]
이다. 따라서 general solution은
\[
    y(x)=c_1e^x+c_2e^{2x}.
\]
Initial conditions를 적용하면
\[
    c_1+c_2=0
\]
이고
\[
    c_1+2c_2=1.
\]
따라서
\[
    c_1=-1,
    \qquad
    c_2=1,
\]
즉
\[
    y(x)=e^{2x}-e^x.
\]
그러므로
\[
    y(\ln2)
    =e^{2\ln2}-e^{\ln2}
    =4-2
    =\boxed{2}.
\]
Constant-coefficient linear ODE에서는 characteristic polynomial을 먼저 세우는 것이 기본이다.
\end{solution}

\begin{problem}{10}
어떤 non-equilateral triangle의 circumcenter와 orthocenter가 각각
\[
    O=(0,0),
    \qquad
    H=(6,3)
\]
이다. 이 triangle의 centroid $G$의 좌표를 구하시오.
\end{problem}

\begin{solution}
Euler Line Theorem에 의해 non-equilateral triangle의 circumcenter $O$, centroid $G$, orthocenter $H$는 collinear이고
\[
    OG:GH=1:2.
\]
따라서 $G$는 $\overline{OH}$를 $1:2$로 internal division하는 점이므로
\[
    G
    =O+\frac13(H-O)
    =\frac13(6,3)
    =(2,1).
\]
따라서
\[
    \boxed{G=(2,1)}.
\]
Euler line에서는 centroid가 circumcenter에서 orthocenter 방향으로 정확히 one third 지점에 놓인다는 형태로 기억하면 계산이 빠르다.
\end{solution}
